\scanpage{294}
\mspage{18}
% archive sheet number: 293

\uline{Proposition} 3.1.
\[
\begin{aligned}
\mathrm{Loc}(6,a)&\Longrightarrow\mathrm{Loc}(6',a)\\
\deleted{\unsure{\mathrm{Loc}(6,b)}}&\deleted{\Longrightarrow\mathrm{Loc}(6'b)\ \&\ \mathrm{Loc}(6'a)}\\
\mathrm{Loc}(6\unsure{,}b)&\Longrightarrow\mathrm{Loc}(6'b)\\
\mathrm{Loc}(6c)&\Longrightarrow\mathrm{Loc}(6'a,b)
\end{aligned}
\]
% slanted
\margin{NB on a même $\mathrm{Loc}\,6c\Longrightarrow\underbrace{\mathrm{Loc}\,6(\deleted{\unsure{c}}a,b)}_{\mathrm{Loc}\,6}$ \deleted{\unsure{$\Longrightarrow$}}}

Ça se voit sur les diagrammes cubiques \deleted{\unsure{S}} $Q$
\diagram{\begin{tikzcd}[ampersand replacement=\&,row sep=small,column sep=small]
X_0\arrow[rr,"i_2"]\arrow[dr,"i_1"'] \& \& X_2\arrow[dr,"j_1"] \&\\
\& X_1\deleted{\unsure{X_1}}\arrow[rr,"{\deleted{\unsure{f_2}}\,j_2}"] \& \& X_3\\
X_0\arrow[uu,"f_0=\mathrm{id}"]\arrow[rr,"i_2",crossing over]\arrow[dr,"\mathrm{id}"'] \& \& X_2\arrow[uu,"{\deleted{\unsure{\mathrm{id}}}f_2=\mathrm{id}}"',crossing over]\arrow[dr,no head,"\mathrm{id}"] \&\\
\& X_0\arrow[uu,"f_1=i_1"]\arrow[rr,"i_2"'] \& \& X_2\arrow[uu,"f_3=j_1"']
\end{tikzcd}\quad $(Q)$,}
\diagram{\begin{tikzcd}[ampersand replacement=\&,row sep=tiny,column sep=tiny]
X_0\arrow[rr,"i_2"]\arrow[dr,"i_1"']\arrow[dd,"i_1"'] \& \& X_2\arrow[dr,"j_1"]\arrow[dd,"j_1",crossing over] \&\\
\&X_1\arrow[rr,"j_2"]\arrow[dd,"\mathrm{id}"'] \& \& X_3\arrow[dd,"\mathrm{id}"]\\
X_1\arrow[rr,"j_2",crossing over]\arrow[dr,"\mathrm{id}"'] \& \& X_3\arrow[dr,"\mathrm{id}"] \&\\
\& X_1\arrow[rr,"j_2"'] \& \& X_3
\end{tikzcd}\quad $(Q')$}
où les $f_0,f_1,f_2,f_3$ de l'énoncé de $\mathrm{Loc}(6)$ deviennent $\mathrm{id}_{X_0},i_1,\mathrm{id}_{X_2},j_1$. $\mathrm{Loc}(6a)$ donne donc $i_1\in W\Longrightarrow j_1\in W$, \inserted{i.e. $\mathrm{Loc}(6'a)$} et $\mathrm{Loc}(6b)$ $j_1\in W\Longrightarrow i_1\in W$, \deleted{\unsure{d'où la}}\inserted{i.e. $\mathrm{Loc}(6'b)$}

\deleted{proposition 3.\unsure{1|2}.} Pour $\mathrm{Loc}(6b)\Longrightarrow\mathrm{Loc}(6',b)$, on utilise \inserted{(cf \uline{XII} p.~\unsure{78|77})} le cube $(Q')$
\diagram{\deleted{\begin{tikzcd}[ampersand replacement=\&,row sep=tiny,column sep=tiny]
\unsure{X_0}\arrow[rr,"{\unsure{i_2}}"]\arrow[dr,"{\unsure{i_1}}"]\arrow[dd,"{\unsure{i_2}}"] \& \& \unsure{X_2}\arrow[dr,"{\unsure{j_1}}"]\arrow[dd,"{\unsure{\mathrm{id}}}"] \&\\
\&\unsure{X_1}\arrow[rr,"{\unsure{j_2}}"]\arrow[dd,"{\unsure{j_2}}"] \& \& X_3\arrow[dd,"{\unsure{\mathrm{id}}}"]\\
\unsure{X_2}\arrow[rr,"{\unsure{\mathrm{id}}}"]\arrow[dr,"{\unsure{j_1}}"] \& \&\unsure{X_2}\arrow[dr,"{\unsure{j_1}}"] \&\\
\& \unsure{X_3}\arrow[rr,"{\unsure{\mathrm{id}}}"] \& \&\unsure{X_3}
\end{tikzcd}}}
, où $f_0=i_1$, $f_1=\mathrm{id}_{X_1}$, $f_{\unsure{2|1}}=j_1$, $f_3=\mathrm{id}_{X_3}$, \inserted{donc $j_1\in W\Longrightarrow i_1\in W$}
% slanted
\margin{\inserted{$\uline{Loc}(6b)\Rightarrow$} i.e. $\uline{Loc}(6'b)$, et si on a $\uline{Loc}(6c)$, on voit que $i_1\in W\Rightarrow j_1\in W$ i.e. $\mathrm{Loc}(6'a)$.}

\uline{Commentaire 3.\unsure{2|?}.} $\mathrm{Loc}(6a)$ m'apparaît comme le ``minimum vital'' pour établir l'existence d'une ``théorie'' des \uline{co}fibres $W$-homotopiques dans $\mathrm{Cat}$. Pour cela il faut en effet disposer d'un critère serviable assurant qu'une flèche $i:A\unsure{-|\longrightarrow}B$ dans $\mathrm{Cat}$ est dans $\mathrm{Cof}_{\unsure{W|w}}$.
